\documentclass[10pt,a4paper]{amsart}
\usepackage{amsmath,amssymb,amsthm,graphicx}
% Retained preamble for the Stacks Project chapter snapshot in this directory.
% The chapter's own preamble.tex is a copy of the shared Stacks Project preamble; the block
% below is sliced MECHANICALLY, line for line, from it. Nothing was added, removed or
% redefined; the wrapper class, the xr-hyper/hyperref wiring and the externaldocument list
% are omitted because the assembler supplies the document class.
%
% lines taken from preamble.tex: 3, 5, 9, 11, 12, 156, 157, 158, ... (51 lines total)
%
% package coverage restored from the chapter's own preamble (2026-10-13): xy, xy-2cell

\usepackage[all]{xy}
\xyoption{2cell}
\UseAllTwocells

\newenvironment{reference}{\comment}{\endcomment}
\newenvironment{slogan}{\comment}{\endcomment}
\newenvironment{history}{\comment}{\endcomment}
\theoremstyle{plain}
\newtheorem{theorem}[subsection]{Theorem}
\newtheorem{proposition}[subsection]{Proposition}
\newtheorem{lemma}[subsection]{Lemma}
\theoremstyle{definition}
\newtheorem{definition}[subsection]{Definition}
\newtheorem{example}[subsection]{Example}
\newtheorem{exercise}[subsection]{Exercise}
\newtheorem{situation}[subsection]{Situation}
\theoremstyle{remark}
\newtheorem{remark}[subsection]{Remark}
\newtheorem{remarks}[subsection]{Remarks}
\numberwithin{equation}{subsection}
\def\lim{\mathop{\mathrm{lim}}\nolimits}
\def\colim{\mathop{\mathrm{colim}}\nolimits}
\def\Spec{\mathop{\mathrm{Spec}}}
\def\Hom{\mathop{\mathrm{Hom}}\nolimits}
\def\Ext{\mathop{\mathrm{Ext}}\nolimits}
\def\SheafHom{\mathop{\mathcal{H}\!\mathit{om}}\nolimits}
\def\SheafExt{\mathop{\mathcal{E}\!\mathit{xt}}\nolimits}
\def\Sch{\mathit{Sch}}
\def\Mor{\mathop{\mathrm{Mor}}\nolimits}
\def\Ob{\mathop{\mathrm{Ob}}\nolimits}
\def\Sh{\mathop{\mathit{Sh}}\nolimits}
\def\NL{\mathop{N\!L}\nolimits}
\def\CH{\mathop{\mathrm{CH}}\nolimits}
\def\proetale{{pro\text{-}\acute{e}tale}}
\def\etale{{\acute{e}tale}}
\def\QCoh{\mathit{QCoh}}
\def\Ker{\mathop{\mathrm{Ker}}}
\def\Im{\mathop{\mathrm{Im}}}
\def\Coker{\mathop{\mathrm{Coker}}}
\def\Coim{\mathop{\mathrm{Coim}}}
\DeclareMathSymbol{\boxtimes}{\mathbin}{AMSa}{"02}
\def\QCohstack{\mathcal{QC}\!\mathit{oh}}
\def\Cohstack{\mathcal{C}\!\mathit{oh}}
\def\Spacesstack{\mathcal{S}\!\mathit{paces}}
\def\Quotfunctor{\mathrm{Quot}}
\def\Hilbfunctor{\mathrm{Hilb}}
\def\Curvesstack{\mathcal{C}\!\mathit{urves}}
\def\Polarizedstack{\mathcal{P}\!\mathit{olarized}}
\def\Complexesstack{\mathcal{C}\!\mathit{omplexes}}
\def\Pic{\mathop{\mathrm{Pic}}\nolimits}
\def\Picardstack{\mathcal{P}\!\mathit{ic}}
\def\Picardfunctor{\mathrm{Pic}}
\def\Deformationcategory{\mathcal{D}\!\mathit{ef}}
\usepackage{cleveref}
\crefname{originaltheorem}{original theorem}{original theorems}
\Crefname{originaltheorem}{Original theorem}{Original theorems}
\makeatletter
\newlabel{injectives-proposition-objects-are-small}{{079F}{source}{}{sourcepointer-1}{}}
\newlabel{injectives-proposition-objects-are-small@cref}{{[originaltheorem][0][]079F}{[1][0][]source}{}{}{}}
\makeatother
\crefname{originaltheorem}{original theorem}{original theorems}
\Crefname{originaltheorem}{Original theorem}{Original theorems}
\makeatletter
\newlabel{lemma-better-set-of-monos}{{0FSY}{source}{}{sourcepointer-2}{}}
\newlabel{lemma-better-set-of-monos@cref}{{[originaltheorem][0][]0FSY}{[1][0][]source}{}{}{}}
\makeatother
\crefname{originaltheorem}{original theorem}{original theorems}
\Crefname{originaltheorem}{Original theorem}{Original theorems}
\makeatletter
\newlabel{lemma-functor-set-of-monos}{{0FSZ}{source}{}{sourcepointer-3}{}}
\newlabel{lemma-functor-set-of-monos@cref}{{[originaltheorem][0][]0FSZ}{[1][0][]source}{}{}{}}
\makeatother
\title{Stacks, Differential Graded Sheaves: theorem-qis-into-dg-injective}
\author{The Stacks Project Authors}
\date{stacks-project:sdga.tex:local-snapshot-20261012; GNU FDL 1.2 or later}
\begin{document}
\maketitle
\noindent\textit{Editorial note.} Statement span is lines 4232-4242 of sources/stacks-sdga/pinned-original-source.tex; proof span is lines 4244-4290, the complete author proof (45 body lines).\par\medskip
\noindent\textit{Editorial note.} Navigation stubs: 3 cited-but-not-included labels are supplied as non-load-bearing pointers with their pinned tags where the snapshot tag index has them.\par\medskip
\noindent\textit{Editorial note.} extra\_packages is empty on purpose: the retained preamble already loads xy with options and a second declaration clashes in the assembler wrapper.\par\medskip
\noindent\textit{Editorial note.} Text is retained verbatim, including the author's own typos and placeholders.\par\medskip
\noindent\textit{Editorial note.} Queue index 9 of queue\_sdga.json (pb=2255, stmt 4232-4242, proof 4244-4290).\par\medskip
\noindent\textit{Allowed standard prerequisites (editorial).}\begin{itemize}
\item Explicit prerequisite boundary. The cited inputs are prerequisites; the argument in the retained span is not replaced by any cited result.
\end{itemize}\medskip
\ifdefined\hypertarget\hypertarget{sourcepointer-1}{}\fi
\noindent\textit{Source reference.} Original theorem 079F: sources/stacks-sdga/pinned-original-source.tex, source lines 4290--4290. Cited by the retained statement or proof but not part of the collected spans; supplied only as a navigation pointer. The argument does not depend on its content.\par\medskip
\ifdefined\hypertarget\hypertarget{sourcepointer-2}{}\fi
\noindent\textit{Source reference.} Original theorem 0FSY: sources/stacks-sdga/pinned-original-source.tex, source lines 4062--4082. Cited by the retained statement or proof but not part of the collected spans; supplied only as a navigation pointer. The argument does not depend on its content.\par\medskip
\ifdefined\hypertarget\hypertarget{sourcepointer-3}{}\fi
\noindent\textit{Source reference.} Original theorem 0FSZ: sources/stacks-sdga/pinned-original-source.tex, source lines 4184--4206. Cited by the retained statement or proof but not part of the collected spans; supplied only as a navigation pointer. The argument does not depend on its content.\par\medskip
\begin{theorem}
\label{theorem-qis-into-dg-injective}
Let $(\mathcal{C}, \mathcal{O})$ be a ringed site.
Let $(\mathcal{A}, \text{d})$
be a sheaf of differential graded algebras on $(\mathcal{C}, \mathcal{O})$.
For every differential graded $\mathcal{A}$-module $\mathcal{M}$ there
exists a quasi-isomorphism $\mathcal{M} \to \mathcal{I}$
where $\mathcal{I}$ is a graded injective and K-injective
differential graded $\mathcal{A}$-module. Moreover, the
construction is functorial in $\mathcal{M}$.
\end{theorem}

\begin{proof}
Let $R$ and $\mathcal{M}_r \to \mathcal{M}'_r$ be a set of
morphisms of $\textit{Mod}(\mathcal{A}, \text{d})$ found in
Lemma \ref{lemma-better-set-of-monos}.
Let $M$ with transformation $\text{id} \to M$
be as constructed in Lemma \ref{lemma-functor-set-of-monos}
using $R$ and $\mathcal{M}_r \to \mathcal{M}'_r$.
Using transfinite recursion we define a sequence of functors
$M_\alpha$ and natural transformations $M_\beta \to M_\alpha$
for $\alpha < \beta$ by setting
\begin{enumerate}
\item $M_0 = \text{id}$,
\item $M_{\alpha + 1} = M \circ M_\alpha$ with natural transformation
$M_\beta \to M_{\alpha + 1}$ for $\beta < \alpha + 1$
coming from the already constructed $M_\beta \to M_\alpha$ and the
maps $M_\alpha \to M \circ M_\alpha$ coming from $\text{id} \to M$, and
\item $M_\alpha = \colim_{\beta < \alpha} M_\beta$ if $\alpha$
is a limit ordinal with the coprojections as transformations
$M_\beta \to M_\alpha$ for $\alpha < \beta$.
\end{enumerate}
Observe that for every differential graded $\mathcal{A}$-module the maps
$\mathcal{M} \to M_\beta(\mathcal{M}) \to M_\alpha(\mathcal{M})$
are injective quasi-isomorphisms (as filtered colimits are exact).

\medskip\noindent
Recall that $\textit{Mod}(\mathcal{A}, \text{d})$ is a Grothendieck
abelian category. Thus by
Injectives, Proposition \ref{injectives-proposition-objects-are-small}
(applied to the direct sum of $\mathcal{M}_r$ for all $r \in R$)
there is a limit ordinal $\alpha$ such that $\mathcal{M}_r$ is $\alpha$-small
with respect to injections for every $r \in R$.
We claim that $\mathcal{M} \to M_\alpha(\mathcal{M})$ is
the desired functorial embedding of $\mathcal{M}$ into a
graded injective K-injective module.

\medskip\noindent
Namely, any map $\mathcal{M}_r \to M_\alpha(\mathcal{M})$
factors through $M_\beta(\mathcal{M})$ for some $\beta < \alpha$.
However, by the construction of $M$ we see that this means
that $\mathcal{M}_r \to M_{\beta + 1}(\mathcal{M}) = M(M_\beta(\mathcal{M}))$
factors through $\mathcal{M}'_r$. Since
$M_\beta(\mathcal{M}) \subset  M_{\beta + 1}(\mathcal{M})
\subset M_\alpha(\mathcal{M})$ we get the desired factorizaton
into $M_\alpha(\mathcal{M})$. We conclude by our choice of
$R$ and $\mathcal{M}_r \to \mathcal{M}'_r$ in
Lemma \ref{lemma-better-set-of-monos}.
\end{proof}

\end{document}
